\documentclass[12pt, a4paper, oneside]{book}

% --- LINGUA E CODIFICA ---
\usepackage[utf8]{inputenc}
\usepackage[T1]{fontenc}
\usepackage[italian]{babel}
\renewcommand{\familydefault}{\sfdefault} % Sans-serif moderno

% --- MARGINI E GEOMETRIA ---
\usepackage[a4paper, margin=2.5cm]{geometry}

% --- COLORI E LINK ---
\usepackage{xcolor}
\definecolor{azzurrino}{HTML}{1E88E5}
\definecolor{grigiochiaro}{HTML}{F5F5F5}
\definecolor{grigioscuro}{HTML}{424242}

\usepackage[
    colorlinks=true,
    linkcolor=azzurrino,
    citecolor=azzurrino,
    urlcolor=azzurrino
]{hyperref}

% --- TITOLI MODERNI ---
\usepackage{titlesec}

\titleformat{\chapter}[display]
  {\normalfont\Huge\bfseries\color{azzurrino}}
  {\filleft\Huge\thechapter}
  {0pt}
  {\titlerule[1pt]\vspace{1ex}\filright}
  [\vspace{1ex}\titlerule]

\titleformat{\section}
  {\normalfont\Large\bfseries\color{azzurrino}}{\thesection}{1em}{}

\titleformat{\subsection}
  {\normalfont\large\bfseries\color{grigioscuro}}{\thesubsection}{1em}{}

\setcounter{tocdepth}{4}
\setcounter{secnumdepth}{4}

\titleformat{\subsubsection}
  {\normalfont\large\bfseries\color{grigioscuro}}{\thesubsubsection}{1em}{}

% --- CAPTION COLORATE ---
\usepackage[font=small, labelfont={bf, color=azzurrino}]{caption}

% --- HEADER E FOOTER ---
\usepackage{fancyhdr}
\setlength{\headheight}{25pt}
\pagestyle{fancy}
\fancyhf{}
\fancyhead[L]{\color{grigioscuro}\leftmark}
\fancyhead[R]{\color{azzurrino}\thepage}
\renewcommand{\headrulewidth}{0.4pt}
\renewcommand{\headrule}{\hbox to\headwidth{\color{azzurrino}\leaders\hrule height \headrulewidth\hfill}}
\fancyfoot[C]{\color{grigioscuro}\small Fondamenti dell'informatica – \thepage}

% --- LISTE ED ENUMERAZIONI ---
\usepackage{enumitem}
\setlist{nosep, leftmargin=1.5em}

% --- GRAFICA E FIGURE ---
\usepackage{graphicx}
\usepackage{bookmark}
\usepackage{float}
\usepackage{wrapfig}
\usepackage[normalem]{ulem} % sottolineature
\usepackage{subcaption}
\usepackage{soul} % evidenziazione
\usepackage{parskip} % spaziatura tra paragrafi

% --- Tcolorbox e Listings ---
\usepackage[most]{tcolorbox}
\tcbuselibrary{listings, skins, breakable, theorems} % theorems utile per ambienti come erexercise
\usepackage{amsmath}
\usepackage{amssymb}

% --- Tabelle ---
\usepackage{longtable}
\usepackage{array}
\usepackage{tabularx}

% --- FRONTESPIZIO ---
\title{\Huge\bfseries Appunti del corso di \\[0.5ex]
\textcolor{azzurrino}{Fondamenti dell'informatica} \\[1ex]
\Large UniVR - Dipartimento di Informatica \\[1.5ex]
\Large Primo semestre - A.A. 2025/26}
\author{\Large Mattia \and Matteo Drago}
\date{}

% --- DEFINIZIONE DI AMBIENTI DI LAVORO ---
%===========================================
% --- AMBIENTE DEFINIZIONE ---
\newcommand{\definizione}[2]{%
  \vspace{5pt}%
  \begin{tcolorbox}[
    colback=azzurrino!5!white,
    colframe=azzurrino!50!black,
    title=\textbf{Definizione: #1},
    coltitle=white,
    fonttitle=\bfseries,
    arc=2mm,
    boxrule=1pt,
    enhanced,
    breakable
  ]
  #2
  \end{tcolorbox}
  \vspace{5pt}%
}

% --- AMBIENTE DEFINIZIONE (Viola) ---
\newcommand{\Definizione}[2]{%
  \vspace{5pt}%
  \begin{tcolorbox}[
    colback=viola!5!white,
    colframe=viola!50!black,
    title=\textbf{Definizione: #1},
    coltitle=white,
    fonttitle=\bfseries,
    arc=2mm,
    boxrule=1pt,
    enhanced,
    breakable
  ]
  #2
  \end{tcolorbox}
  \vspace{5pt}%
}

% --- AMBIENTE DEFINIZIONE (Verde) ---
\newcommand{\definition}[2]{%
  \vspace{5pt}%
  \begin{tcolorbox}[
    colback=verde!5!white,
    colframe=verde!50!black,
    title=\textbf{Definizione: #1},
    coltitle=white,
    fonttitle=\bfseries,
    arc=2mm,
    boxrule=1pt,
    enhanced,
    breakable
  ]
  #2
  \end{tcolorbox}
  \vspace{5pt}%
}

%===========================================
% --- AMBIENTE ESEMPIO ---
\newcommand{\esempio}[2]{%
  \vspace{5pt}%
  \begin{tcolorbox}[
    colback=brown!5!white,
    colframe=brown!50!black,
    title=\textbf{Esempio: #1},
    coltitle=white,
    fonttitle=\bfseries,
    arc=2mm,
    boxrule=1pt,
    enhanced,
    breakable
  ]
  #2
  \end{tcolorbox}
  \vspace{5pt}%
}

%===========================================
% --- AMBIENTE VARIO ---
\newcommand{\vario}[2]{%
  \vspace{5pt}%
  \begin{tcolorbox}[
    colback=orange!10!white,
    colframe=orange!70!black,
    title=\textbf{#1},
    coltitle=white,
    fonttitle=\bfseries,
    arc=2mm,
    boxrule=1pt,
    enhanced,
    breakable
  ]
  #2
  \end{tcolorbox}
  \vspace{5pt}%
}

%===========================================
% --- STILI PER I LINGUAGGI DI PROGRAMMAZIONE ---
%===========================================
\lstdefinestyle{Rstyle}{
  language=R,
  basicstyle=\ttfamily\small,
  keywordstyle=\color{blue},
  stringstyle=\color{red},
  keepspaces=true,       % <-- Rispetta gli spazi
  columns=fixed,  % <-- Impedisce a LaTeX di allargare il testo
  commentstyle=\color{green!50!black}, % NESSUN CANCELLETTO QUI, ERRORE RISOLTO!
  morekeywords={c, length, print, summary, mean, sd, var, min, max, sum, 
                seq, rep, paste, paste0, rbind, cbind, matrix, data.frame, 
                list, str, class, typeof, apply, lapply, sapply, tapply, 
                plot, hist, boxplot, read.csv, write.csv, head, tail, 
                subset, merge, table, ifelse, require, install.packages, library},
  numbers=left,
  numberstyle=\tiny\color{gray},
  breaklines=true,
  showstringspaces=false
}

\lstdefinestyle{SQLstyle}{
  language=SQL,
  basicstyle=\ttfamily\footnotesize,
  keywordstyle=\color{keywordcolor}\bfseries,
  commentstyle=\color{gray},
  stringstyle=\color{red},
  breaklines=true,
  showstringspaces=false,
  morekeywords={DISTINCT, AS, FROM, WHERE, GROUP, BY, HAVING, UNION, INTERSECT, EXCEPT, ORDER, ASC, DESC, USING}
}

\lstdefinestyle{Javastyle}{
  language=Java,
  basicstyle=\ttfamily\footnotesize,
  keywordstyle=\color{viola}\bfseries,
  stringstyle=\color{red},
  commentstyle=\color{green!60!black}\itshape,
  breaklines=true,
  showstringspaces=false
}

%===========================================
% --- AMBIENTI CODICE (SENZA TITOLO) ---
%===========================================
\newtcblisting{codiceR}[1][]{%
  width=\linewidth, % <--- AGGIUNGI QUESTA RIGA!
  colback=yellow!10!white,
  colframe=yellow!50!black,
  sharp corners, % <-- Angoli vivi, zero arrotondamenti
  boxrule=0.8pt,
  top=0pt,       % <-- Rimuove la riga vuota interna superiore
  bottom=0pt,    % <-- Rimuove la riga vuota interna inferiore
  left=2pt,      % Un minimo indispensabile per non appiccicare le lettere al bordo
  right=2pt,
  enhanced,
  breakable,
  listing only,
  listing options={style=Rstyle}, 
  #1 
}

\newtcblisting{codiceSQL}[1][]{%
  colback=blue!3!white,     % Azzurrino chiarissimo
  colframe=blue!40!black,   % Bordo blu neutro
  arc=2mm,
  boxrule=0.8pt,
  enhanced,
  breakable,
  listing only,
  listing options={style=SQLstyle},
  #1
}

\newtcblisting{codiceJava}[1][]{%
  colback=orange!5!white,   % Arancione molto tenue
  colframe=orange!60!black, % Bordo arancione bruciato
  arc=2mm,
  boxrule=0.8pt,
  enhanced,
  breakable,
  listing only,
  listing options={style=Javastyle},
  #1
}

%==========================================
% --- AMBIENTI ER E MATEMATICI ---
%==========================================
\newtcolorbox{er}[2][]{%
  enhanced, breakable,
  colback=gray!3,
  colframe=blue!50!black,
  fonttitle=\bfseries,
  title={Esercizio E-R: #2},
  before upper={\refstepcounter{erex}},
  sharp corners=south,
  boxrule=0.8pt,
  arc = 2mm,
  #1
}

\newcommand{\teorema}[2]{%
  \vspace{5pt}%
  \begin{tcolorbox}[
    colback=red!5!white,
    colframe=red!50!black,
    title=\textbf{Teorema: #1},
    coltitle=white,
    fonttitle=\bfseries,
    arc=2mm,
    boxrule=1pt,
    enhanced,
    breakable
  ]
  #2
  \end{tcolorbox}
  \vspace{5pt}%
}

\newcommand{\lemma}[3]{%
  \vspace{5pt}%
  \begin{tcolorbox}[
    colback=red!5!white,
    colframe=red!50!black,
    title=\textbf{Lemma: #1},
    coltitle=white,
    fonttitle=\bfseries,
    arc=2mm,
    boxrule=1pt,
    enhanced,
    breakable
  ]
  #2
  \ifstrempty{#3}{}{%
    \par\vspace{8pt}%
    \textbf{Dimostrazione}\\
    #3
  }
  \end{tcolorbox}
  \vspace{5pt}%
}

\newcommand{\osservazione}[1]{%
  \vspace{5pt}%
  \begin{tcolorbox}[
    colback=white!5!white,
    colframe=white!50!black,
    title=\textbf{Osservazione},
    coltitle=white,
    fonttitle=\bfseries,
    arc=2mm,
    boxrule=1pt,
    enhanced,
    breakable
  ]
  #1
  \end{tcolorbox}
  \vspace{5pt}%
}

\newtcolorbox{esempiobox}[1]{
    colback=brown!5, 
    colframe=brown, 
    title=#1, 
    fonttitle=\bfseries
} %file apposito per definire nuovi comandi

% =================================
\begin{document}
    \frontmatter
    \maketitle
    \tableofcontents
    \thispagestyle{empty}
    \newpage

    %===========================
    \mainmatter

    % --- CAPITOLO 1: NOZIONI DI BASE DEL CORSO ---
    \chapter{Nozioni di base}

    \section{Cardinalità degli insiemi}

    \teorema{Ciao}{Ciao}





    % --- CAPITOLO 2: Linguaggi formali ---
    \chapter{Linguaggi formali}

    \section{Linguaggi regolari}
    I \textbf{linguaggi regolari} rappresentano la famiglia dei modelli di calcolo più semplice: \hl{corrispondono al processore di un calcolatore}.

    I problemi relativi a questi linguaggi sono \textit{decisionali}: una sequenza di simboli (stringa) appartiene a un linguaggio regolare, allora è risolvibile da un algoritmo e implementabile mediante una macchina stati finiti.

    \definizione{stringhe}{Un \textbf{simbolo} è un'entità primitiva astratta non definita formalmente (lettere e caratteri numerici).}

    \definizione{alfabeto}{Un \textbf{alfabeto} è un'insieme di simboli denotato con $\Sigma$ e avente le seguenti caratteristiche:
    \begin{itemize}
      \item \textbf{cardinalità di $\Sigma$ minore di $\omega$}:
      \begin{equation*}
        |\Sigma| < \omega
      \end{equation*}
      \item \textbf{insieme $\Sigma \neq 0$} $\Rightarrow$ è sempre presente almeno un simbolo
    \end{itemize}
    }

    L'\hl{\textbf{insieme delle stringhe} sull'alfabeto $\Sigma$ è rappresentato con $\Sigma$*}: insieme di tutte le sequenze di simboli che appartegono all'alfabeto $\Sigma$.

    \esempio{insieme delle stringhe $\Sigma$*}{Se $\Sigma$ è rappresentato dall'insieme binario $\{0, 1\}$, allora la stringa 101011 appartiene a $\{0, 1\}$*:
    \begin{equation*}
      \Sigma = \{0, 1\}, \ \ \ 101011 \in \{0, 1\}^*
    \end{equation*}}

    La \textit{cardinalità} dell'insieme delle stringhe $\Sigma$*, su qualunque alfabeto, è uguale alla cardinalità dei numeri naturali:
   \begin{equation*}
    |\Sigma^*| = |\mathbb{N}| = \omega
  \end{equation*}

  A partire dagli insiemi di sequenze di simboli (stringhe) si formano i \textbf{linguaggi}.

  \definizione{linguaggio}{Un linguaggio $\mathcal{L}$ è un sottoinisieme dell'insieme delle stringhe:
  \begin{equation*}
    \mathcal{L} \subseteq \Sigma^*
  \end{equation*}}

  \textit{Quanti sono i linguaggi?} \\[0.2ex]
  Se la cardinalità di $\Sigma$* è pari alla cardinalità dei numeri naturali, allora la cardinalità dei sottoinisiemi dei numer naturali è ugale alla Cardinalità dei sottoinisiemi degli insiemi delle stringhe:
  \begin{equation*}
    |\Sigma^*| = |\mathbb{N}| = \omega 
  \end{equation*}
  \begin{center}
    $\Downarrow$
  \end{center}
  \begin{equation*}
    |2^N| = |2^{\Sigma^*}| = |\mathbb{N} \rightarrow |\mathbb{N}| = |\mathbb{R}|
  \end{equation*}

  Di conseguenza ci sono tanti linguaggi quanti i numeri reali su un alfabeto finito:
  \begin{equation*}
    |\{ \mathcal{L} | \mathcal{L} \subseteq \Sigma^*\}| \subseteq |\mathbb{R}|
  \end{equation*}




  \subsection{Autommi a stati finiti}
  Un \textbf{automa a stati finiti} è un dispositivo avente un nastro in \textit{input}, suddiviso a blocchi, su cui la macchina non può scrivere, bensì solo leggere.

  %foto

  La stringa $\sigma$ è formata dai simboli $x_0, x_1, x_2, \dots$, mentre il \textbf{controllore} (dipositivo di lettura) parte da un certo stato $q$ iniziale e può muoversi soltanto da sinistra verso destra. \\[0.2ex]
  Quindi si può leggere un simbolo $x_i$, consumare questo simbolo, effetture una transizione di stato passando da $q$ a $q^t$ e, infine, passare al simbolo sucessivo $x_{i+1}$. \\[0.2ex]
  Man mano che la testina procede, essa consuma ciò che è presente sul nastro in \textit{input} fintatochè la stringa $\sigma$ non si esaurisce.

  Poichè le stringhe che si possono mettere sul nastro sono finite, allora le operazioni indicate possono continuare fino a un certo numero di passi, dopodichè la macchina termina.

  Una descrizione matematicaa di questa $M$ consiste nella quintupla:
  \begin{equation*}
    M = \langle Q, \Sigma, \delta, q_0, F\rangle
  \end{equation*}

  \begin{description}
    \item [$Q$]: insieme di \textit{stati}
    \item [$\Sigma$]: insieme di \textit{simboli}
    \item [$\delta$]: funzione di \textit{transizione} 
    \item [$q_0$]: stato inizale da cui la macchina parte
    \item [$F$]: insieme di stati finali
  \end{description}

  \esempio{automa a stato finito}{
    %disegno
    \begin{itemize}
      \item singola freccia entrante $\rightarrow$ \textit{stato iniziale}
      \item doppio cercio $\rightarrow$ \textit{stato finale}
      \item frecce $\rightarrow$ \textit{archi} che a seconda del valore letto effettuano una transizione dal nodo cambiando lo stato
    \end{itemize}
  }

  Il linguaggio accettato dall'automa è l'insieme di tutte òe stringhe sull'alfabeto tale che $\hat{\delta}(q_0, \sigma) \in F$:
  \begin{equation*}
    \mathcal{L}(M) = \{\sigma \in \Sigma^* | \hat{\delta}(q_0, \sigma) \in F\}
  \end{equation*}

  \definizione{linguaggio regolare}{Se $M$ è un'automa a stati finiti deterministico, allora $\mathcal{L}$ è detto \textbf{regolare}.}

  \subsubsection{Automa a stati finiti deterministico}



\end{document}